\documentclass[11pt,a4paper]{article}

% ── Packages ──────────────────────────────────────────────────────────────
\usepackage[utf8]{inputenc}
\usepackage[T1]{fontenc}
\usepackage{amsmath,amssymb,amsthm}
\usepackage{mathtools}
\usepackage{bm}
\usepackage{geometry}
\geometry{margin=1in,headheight=14pt}
\usepackage{hyperref}
\hypersetup{colorlinks=true,linkcolor=blue!70!black,citecolor=blue!70!black,urlcolor=blue!70!black}
\usepackage{xcolor}
\usepackage{titlesec}
\usepackage{enumitem}
\usepackage{booktabs}
\usepackage{tabularx}
\usepackage{multirow}
\usepackage{fancyhdr}
\usepackage{graphicx}
\usepackage{float}
\usepackage{caption}
\usepackage{subcaption}
\usepackage{algorithm}
\usepackage{algpseudocode}

% ── Theorem environments ─────────────────────────────────────────────────
\newtheorem{definition}{Definition}[section]
\newtheorem{theorem}{Theorem}[section]
\newtheorem{remark}{Remark}[section]

% ── Macros ────────────────────────────────────────────────────────────────
\newcommand{\R}{\mathbb{R}}
\newcommand{\Hk}{H_k}
\newcommand{\Hkp}{H_{k+1}}
\newcommand{\sk}{s_k}
\newcommand{\yk}{y_k}
\newcommand{\gk}{g_k}
\newcommand{\xk}{x_k}
\newcommand{\tp}{\tau_k}
\newcommand{\thk}{\theta_k}
\newcommand{\phik}{\varphi_k}
\DeclareMathOperator{\tr}{tr}
\DeclareMathOperator{\diag}{diag}
\DeclareMathOperator{\lap}{\nabla^2}

% ── Header / Footer ──────────────────────────────────────────────────────
\pagestyle{fancy}
\fancyhf{}
\rhead{Cahn--Hilliard PINN Experiments}
\lhead{DRP --- Optimizing the Optimizer}
\cfoot{\thepage}

% ── Title ─────────────────────────────────────────────────────────────────
\title{%
  \textbf{Optimizing Quasi-Newton Methods for\\
  Physics-Informed Neural Networks}\\[6pt]
  \large A Comparative Study on the 2D Cahn--Hilliard Equation\\[4pt]
  \normalsize DRP Experiment Report
}
\author{Maria Maciaso-Lugo}
\date{February 2026}

% ══════════════════════════════════════════════════════════════════════════
\begin{document}
\maketitle
\tableofcontents
\newpage

% ======================================================================
%  CHAPTER 1 — INTRODUCTION
% ======================================================================
\section{Introduction}
\label{sec:intro}

Physics-Informed Neural Networks (PINNs) embed the governing partial differential
equations (PDEs) directly into the loss function of a neural network, enabling
mesh-free solution of forward and inverse problems.  A critical bottleneck in
PINN training is the \emph{optimizer}: first-order methods like Adam converge
quickly to the vicinity of a minimum but stall at moderate accuracy, while
second-order methods like BFGS can achieve much higher precision if properly
configured.

This report documents a systematic ablation study of \textbf{17 experiments}
on the 2D Cahn--Hilliard equation, testing:
\begin{itemize}[nosep]
  \item Six quasi-Newton optimizer variants (BFGS, SSBFGS-OL, SSBFGS-AB, SSBroyden1, SSBroyden2, L-BFGS-B)
  \item Adam warmup duration (0, 500, 2000, 5000 epochs)
  \item Adam learning rate (0.0001, 0.001, 0.01)
  \item Collocation resampling frequency (every 200, 500, 1000 iterations)
  \item Power transforms for Hessian conditioning ($p = 1, 2, 4$)
  \item Nocedal--Wright initial Hessian scaling
  \item Adam-only baseline (12,000 epochs)
\end{itemize}

All experiments share a fixed total budget of \textbf{12,000 iterations},
a network with \textbf{1,361 trainable parameters}, and identical random
seeds, sampling strategies, and loss weights.  The only variable in each
experiment is the specific hyperparameter under study.

The goal is to identify which optimizer configuration yields the lowest
final loss, best mass conservation, and most favorable loss-vs-compute
tradeoff for the Cahn--Hilliard PINN.


% ======================================================================
%  CHAPTER 2 — THE CAHN--HILLIARD EQUATION
% ======================================================================
\section{The Cahn--Hilliard Equation}
\label{sec:ch}

\subsection{Physical Background}

The Cahn--Hilliard equation models \textbf{spinodal decomposition} and
phase separation in binary mixtures.  The order parameter $U(t, x, y)
\in [-1, 1]$ represents the local concentration difference between two
phases.  Starting from a nearly uniform mixture, the system spontaneously
separates into domains of $U \approx +1$ and $U \approx -1$, connected
by smooth interfaces whose width is controlled by a diffusivity parameter.

\subsection{The PDE}

Following the Model~B dynamics formulation, the Cahn--Hilliard equation
in two spatial dimensions is:
\begin{equation}
\boxed{
  \frac{\partial U}{\partial t}
  = D \,\nabla^{2}\!\Big(
      \nabla^{2} U + a_2\, U + a_4\, U^3
    \Big)
}
\label{eq:ch}
\end{equation}
where:
\begin{itemize}[nosep]
  \item $D > 0$ is the diffusion coefficient (we use $D = 1$),
  \item $a_2 = 1$ and $a_4 = 1$ are parameters of the double-well free energy
        $f(U) = -\tfrac{a_2}{2} U^2 + \tfrac{a_4}{4} U^4$,
  \item $\nabla^2 = \partial_{xx} + \partial_{yy}$ is the Laplacian,
  \item The domain is $\Omega = [0, 128]^2$ with periodic boundary conditions.
\end{itemize}

\subsection{Expanding the Operator}

The term $\nabla^2(U^3)$ requires careful expansion via the chain rule.
Let $U^3$ be a function of $(x,y)$ through $U(t,x,y)$:
\begin{align}
  \nabla^2(U^3)
  &= \frac{\partial^2}{\partial x^2}(U^3) + \frac{\partial^2}{\partial y^2}(U^3) \notag\\
  &= 6U\left(U_x^2 + U_y^2\right) + 3U^2\left(U_{xx} + U_{yy}\right) \notag\\
  &= 6U\,|\nabla U|^2 + 3U^2\,\nabla^2 U.
  \label{eq:lap_U3}
\end{align}
Substituting back and expanding $\nabla^2(\nabla^2 U) = \nabla^4 U$
(the biharmonic operator), the full PDE residual is:
\begin{equation}
  r(t,x,y) = U_t - D\Big[
    \underbrace{U_{xxxx} + 2\,U_{xxyy} + U_{yyyy}}_{\nabla^4 U}
    + a_2\underbrace{(U_{xx} + U_{yy})}_{\nabla^2 U}
    + a_4\big(6U\,(U_x^2 + U_y^2) + 3U^2\,(U_{xx} + U_{yy})\big)
  \Big].
  \label{eq:residual}
\end{equation}

\subsection{Conservation Law}

A key physical property of the Cahn--Hilliard equation with periodic
boundary conditions is \textbf{mass conservation}:
\begin{equation}
  \frac{d}{dt} \int_\Omega U \, dx \, dy = 0.
  \label{eq:mass}
\end{equation}
This follows because the right-hand side of \eqref{eq:ch} is a divergence
(Laplacian of a quantity), and its integral vanishes under periodic BCs.
We use mass drift $\Delta m = \max_t \langle U \rangle_\Omega - \min_t \langle U \rangle_\Omega$
as a diagnostic metric.

\subsection{Domain and Parameters}

\begin{center}
\begin{tabular}{lll}
\toprule
\textbf{Parameter} & \textbf{Value} & \textbf{Description} \\
\midrule
$(x, y) \in \Omega$  & $[0, 128]^2$ & Spatial domain \\
$t \in [0, T]$       & $[0, 20]$    & Time window \\
$D$                   & 1.0          & Diffusion coefficient \\
$a_2, a_4$            & 1.0, 1.0     & Free energy parameters \\
BCs                   & Periodic     & Soft enforcement (penalty) \\
IC                    & $U_0 \sim \text{Uniform}(-1, 1)$, smoothed & Gaussian-smoothed random field \\
\bottomrule
\end{tabular}
\end{center}


% ======================================================================
%  CHAPTER 3 — PINN FORMULATION
% ======================================================================
\section{PINN Formulation}
\label{sec:pinn}

\subsection{Network Architecture}

The PINN approximates the solution $U(t,x,y)$ with a fully connected
neural network $\hat{U}_\theta : \R^3 \to \R$:
\begin{equation}
  \hat{U}_\theta(t, x, y) = W_L \,\sigma\!\big(W_{L-1}\,\sigma(\cdots\sigma(W_1 z_0 + b_1)\cdots) + b_{L-1}\big) + b_L
\end{equation}
where $z_0 = [t, x, y]^\top$ is the input (after normalization to $[-1,1]^3$),
$\sigma(\cdot) = \tanh(\cdot)$, and the layer dimensions are $[3, 20, 20, 20, 20, 1]$.
This yields:
\begin{equation}
  n_{\text{params}} = (3 \times 20 + 20) + (20 \times 20 + 20) \times 3 + (20 \times 1 + 1) = 1{,}361.
\end{equation}

\subsection{Loss Function}

The training objective combines three terms:
\begin{equation}
\boxed{
  \mathcal{L}(\theta) = \lambda_{\text{pde}}\,\mathcal{L}_{\text{pde}}
                       + \lambda_{\text{ic}}\,\mathcal{L}_{\text{ic}}
                       + \lambda_{\text{bc}}\,\mathcal{L}_{\text{bc}}
}
\label{eq:loss}
\end{equation}
\begin{align}
  \mathcal{L}_{\text{pde}} &= \frac{1}{N_{\text{int}}} \sum_{i=1}^{N_{\text{int}}}
    \big| r(t_i, x_i, y_i) \big|^2,
  \label{eq:loss_pde} \\
  \mathcal{L}_{\text{ic}}  &= \frac{1}{N_0} \sum_{i=1}^{N_0}
    \big| \hat{U}_\theta(0, x_i, y_i) - U_0(x_i, y_i) \big|^2,
  \label{eq:loss_ic} \\
  \mathcal{L}_{\text{bc}}  &= \text{MSE of periodic matching}
    \;\big(\hat{U}_\theta \text{ and } \partial_n \hat{U}_\theta
    \text{ on opposing faces}\big).
  \label{eq:loss_bc}
\end{align}
We use $\lambda_{\text{pde}} = 1$, $\lambda_{\text{ic}} = 5$,
$\lambda_{\text{bc}} = 5$, emphasizing boundary and initial conditions.

\subsection{Collocation Sampling and RAD}

Points are sampled randomly:
$N_{\text{int}} = 10{,}000$ interior, $N_0 = 1{,}000$ initial,
$N_b = 500$ boundary (per face pair).  \textbf{Residual-Adaptive
Distribution (RAD)} resamples interior points with probability
proportional to $|r|^{k_1}$ (with $k_1 = k_2 = 1$), concentrating
collocation points where the PDE residual is large.  Resampling
occurs periodically --- the frequency is one of our ablation axes.

\subsection{Two-Phase Training}

The canonical training pipeline follows a two-phase strategy:
\begin{enumerate}[nosep]
  \item \textbf{Phase 1 (Adam):} First-order optimizer for global
        exploration.  Decaying learning rate:
        $\text{lr}(t) = \text{lr}_0 \cdot 0.98^{t/1000}$.
  \item \textbf{Phase 2 (Quasi-Newton):} SciPy's \texttt{minimize}
        with a patched BFGS implementation supporting multiple
        self-scaling variants.  The inverse Hessian approximation $H_k$
        is \textbf{warm-started} across resampling events.
\end{enumerate}
The total budget is fixed at 12,000 iterations across both phases.


% ======================================================================
%  CHAPTER 4 — THE MATH OF QUASI-NEWTON METHODS
% ======================================================================
\section{Quasi-Newton Optimizer Variants: The Mathematics}
\label{sec:optimizers}

All variants tested share the same outer loop: given a current iterate
$\xk$, gradient $\gk = \nabla f(\xk)$, and inverse Hessian approximation
$\Hk \approx [\nabla^2 f(\xk)]^{-1}$, compute:
\begin{enumerate}[nosep]
  \item \textbf{Search direction:} $p_k = -\Hk \gk$
  \item \textbf{Line search:} find $\alpha_k > 0$ satisfying Wolfe conditions
  \item \textbf{Update:} $x_{k+1} = \xk + \alpha_k p_k$
  \item \textbf{Curvature pairs:} $\sk = \alpha_k p_k$, $\yk = g_{k+1} - \gk$
  \item \textbf{Hessian update:} $\Hkp = \text{Update}(\Hk, \sk, \yk)$
        \quad (this is where variants differ)
\end{enumerate}

We define the following key quantities used across all update formulae:
\begin{align}
  \rho_k &= \frac{1}{\yk^\top \sk}, \label{eq:rhok}\\
  h_k &= \rho_k \,\yk^\top \Hk \yk, \label{eq:hk}\\
  b_k &= -\alpha_k \,\rho_k \,\sk^\top\!(g_{k+1} - \yk), \label{eq:bk}\\
  a_k &= b_k h_k - 1. \label{eq:ak}
\end{align}
The quantity $\rho_k$ is the inverse curvature, $h_k$ measures how well
$\Hk$ captures the curvature along $\yk$, and $b_k$ encodes information
about the line search quality.

% ------------------------------------------------------------------
\subsection{Standard BFGS}
\label{sec:bfgs}

The classical BFGS update maintains positive definiteness of $\Hk$ as long
as $\yk^\top \sk > 0$ (guaranteed by the Wolfe line search):

\begin{equation}
\boxed{
  \Hkp = \Hk
    - \rho_k\!\big(\Hk \yk \,\sk^\top + \sk\, \yk^\top \Hk\big)
    + \rho_k(1 + h_k)\,\sk \sk^\top
}
\label{eq:bfgs}
\end{equation}

This is equivalent to the well-known Sherman--Morrison--Woodbury form:
\begin{equation}
  \Hkp = (I - \rho_k \sk \yk^\top)\,\Hk\,(I - \rho_k \yk \sk^\top)
          + \rho_k \sk \sk^\top.
  \label{eq:bfgs_smw}
\end{equation}
\textbf{Key property:} BFGS satisfies the \emph{secant condition}
$\Hkp \yk = \sk$, ensuring the update stays consistent with observed
curvature.

\textbf{Storage:} $O(n^2)$ for the full $\Hk$ matrix (with
$n = 1{,}361$, this is about $15$ MB in float64).

% ------------------------------------------------------------------
\subsection{L-BFGS-B (Limited-Memory BFGS)}
\label{sec:lbfgs}

L-BFGS-B avoids storing the full $n \times n$ matrix by keeping only
the $m$ most recent curvature pairs $\{(\sk, \yk)\}_{k-m+1}^{k}$
(we use $m = 20$).  The inverse Hessian--vector product $\Hk \gk$ is
computed implicitly via the \textbf{two-loop recursion}:

\begin{algorithm}[H]
\caption{L-BFGS Two-Loop Recursion}
\begin{algorithmic}[1]
\State $q \gets \gk$
\For{$i = k, k{-}1, \ldots, k{-}m{+}1$}
  \State $\alpha_i \gets \rho_i \, s_i^\top q$
  \State $q \gets q - \alpha_i \, y_i$
\EndFor
\State $r \gets \gamma_k \cdot q$ \Comment{$\gamma_k = s_k^\top y_k / y_k^\top y_k$}
\For{$i = k{-}m{+}1, \ldots, k$}
  \State $\beta \gets \rho_i \, y_i^\top r$
  \State $r \gets r + (\alpha_i - \beta)\, s_i$
\EndFor
\State \Return $r = \Hk \gk$
\end{algorithmic}
\end{algorithm}

\textbf{Storage:} $O(mn)$ instead of $O(n^2)$.

\textbf{Tradeoffs:}
\begin{itemize}[nosep]
  \item[\textcolor{green!60!black}{$+$}] Memory-efficient: only $2mn$ floats ($\approx 0.4$ MB vs 15 MB).
  \item[\textcolor{green!60!black}{$+$}] No Cholesky check needed.
  \item[\textcolor{red!70!black}{$-$}] Loses curvature information beyond the last $m$ pairs.
  \item[\textcolor{red!70!black}{$-$}] Cannot warm-start $\Hk$ across collocation resampling events.
\end{itemize}

% ------------------------------------------------------------------
\subsection{Self-Scaled BFGS Methods}
\label{sec:ssbfgs}

Self-scaling modifies the BFGS update by multiplying $\Hk$ by a
\textbf{scaling factor} $\tp > 0$ before the rank-two correction.
The general self-scaled BFGS update is:
\begin{equation}
\boxed{
  \Hkp = \frac{1}{\tp}\!\left[\Hk
    - \rho_k\big(\Hk \yk\, \sk^\top + \sk\, \yk^\top \Hk\big)
    + \rho_k(1 + h_k)\,\sk \sk^\top\right]
    + \rho_k \,\sk \sk^\top \left(1 - \frac{1}{\tp}\right)
}
\label{eq:ssbfgs_general}
\end{equation}
which simplifies to the identity when $\tp = 1$ (standard BFGS).
The choice of $\tp$ distinguishes the variants.

% ------------------------------------------------------------------
\subsubsection{SSBFGS-OL: Oren--Luenberger Scaling}
\label{sec:ssbfgs_ol}

Oren and Luenberger (1974) proposed choosing $\tp$ to minimize the
condition number of $\Hkp$ along a one-parameter family:
\begin{equation}
\boxed{
  \tp^{\text{OL}} = \frac{1}{b_k}
  = \frac{-1}{\alpha_k \rho_k \, \sk^\top(g_{k+1} - \yk)}
}
\label{eq:tau_ol}
\end{equation}
This makes the scaling $\Hk \gets \Hk / \tp$ before the BFGS update.
When $b_k > 1$, the inverse Hessian is shrunk; when $b_k < 1$, it is expanded.

\begin{remark}
The factor $b_k$ can be interpreted as a measure of how well the
current $\Hk$ predicts the actual curvature change.  When $b_k = 1$,
the scaling is neutral and the update reduces to standard BFGS.
\end{remark}

% ------------------------------------------------------------------
\subsubsection{SSBFGS-AB: Al-Baali Scaling}
\label{sec:ssbfgs_ab}

Al-Baali (1998) proposed a more conservative safeguard:
\begin{equation}
\boxed{
  \tp^{\text{AB}} = \min\!\left(1, \;\frac{1}{b_k}\right)
}
\label{eq:tau_ab}
\end{equation}
This ensures $\tp \leq 1$, so the inverse Hessian is \emph{never shrunk}---only expanded or left unchanged.  The intuition is that shrinking $\Hk$
(equivalently, inflating the effective Hessian $B_k = \Hk^{-1}$) can lead to
overly conservative steps, while expanding it can only make steps larger.

\begin{remark}
Compare with OL: when $b_k > 1$, OL scales by $1/b_k < 1$ (shrinks $\Hk$),
but AB clamps at $\tp = 1$ (no scaling).  When $b_k < 1$, both agree:
$\tp = 1/b_k > 1$ (expands $\Hk$).
\end{remark}

% ------------------------------------------------------------------
\subsection{Self-Scaled Broyden Class}
\label{sec:ssbroyden}

The Broyden class generalizes both BFGS ($\thk = 0$) and DFP ($\thk = 1$)
through a convex combination parameter $\thk \in [0, 1]$:
\begin{equation}
  \Hkp^{\text{Broyden}} = \Hkp^{\text{BFGS}} + \thk\,(\yk^\top \Hk \yk)\,v_k v_k^\top
  \label{eq:broyden_class}
\end{equation}
where $v_k = \rho_k \sk - \Hk \yk / (\yk^\top \Hk \yk)$.

The \textbf{self-scaled} version combines this with both $\tp$ and $\thk$:
\begin{equation}
\boxed{
  \Hkp = \frac{1}{\tp}\left[
    \Hk - \frac{\Hk \yk\, \yk^\top \Hk}{\yk^\top \Hk \yk}
    + \phik \,(\yk^\top \Hk \yk)\, v_k v_k^\top
  \right] + \rho_k \,\sk \sk^\top
}
\label{eq:ssbroyden}
\end{equation}
where:
\begin{equation}
  \phik = \frac{1 - \thk}{1 + a_k\,\thk}, \qquad
  \sigma_k = 1 + \thk a_k.
  \label{eq:phik}
\end{equation}

% ------------------------------------------------------------------
\subsubsection{SSBroyden1: DFP Direction ($\thk$: max toward DFP)}
\label{sec:ssbroyden1}

SSBroyden1 chooses $\thk$ to move toward the DFP end of the Broyden
family.  The algorithm first computes:
\begin{align}
  \rho_k^{(-)} &= \min\!\Big(1,\; h_k\Big(1 - \sqrt{\tfrac{|a_k|}{1+a_k}}\,\Big)\Big),
  \label{eq:rhokm}\\
  \thk^{(-)} &= \frac{\rho_k^{(-)} - 1}{a_k}, \qquad
  \thk^{(+)} = \frac{1}{\rho_k^{(-)}}, \qquad
  \thk^{\text{SR1}} = \frac{1}{1 - b_k}.
  \label{eq:theta_candidates}
\end{align}
Then $\thk$ is selected by a case switch on $b_k$:
\begin{equation}
  \thk = \begin{cases}
    \thk^{(-)} & \text{if } b_k = 1, \\
    \max(\thk^{(-)},\; \thk^{\text{SR1}}) & \text{if } b_k > 1, \\
    \min(\thk^{(+)},\; \thk^{\text{SR1}}) & \text{if } b_k < 1.
  \end{cases}
  \label{eq:theta_broyden1}
\end{equation}
The scaling $\tp$ is then determined by $\sigma_k = 1 + \thk a_k$ and
a dimension-dependent term:
\begin{equation}
  \sigma_k^{1/(1-n)} = |\sigma_k|^{1/(1-n)}, \qquad
  \tp = \begin{cases}
    \min\!\big(\min(1, 1/b_k) \cdot \sigma_k^{1/(1-n)},\; \sigma_k\big) & \text{if } \thk \leq 0, \\
    \min(1, 1/b_k) \cdot \min\!\big(\sigma_k^{1/(1-n)},\; 1/\thk\big) & \text{if } \thk > 0.
  \end{cases}
  \label{eq:tau_broyden1}
\end{equation}

% ------------------------------------------------------------------
\subsubsection{SSBroyden2: Optimized $\thk$}
\label{sec:ssbroyden2}

SSBroyden2 replaces the case-switch logic with a simpler, more robust
clamped formula:
\begin{equation}
\boxed{
  \thk = \max\!\left(\thk^{(-)},\;\min\!\left(\thk^{(+)},\;\frac{1-b_k}{b_k}\right)\right)
}
\label{eq:theta_broyden2}
\end{equation}
This clamps $\thk$ to $[\thk^{(-)}, \thk^{(+)}]$ with a target of
$(1-b_k)/b_k$ that interpolates smoothly between BFGS-like and DFP-like
behavior based on the line search quality $b_k$.

The scaling formula for $\tp$ is identical to SSBroyden1 \eqref{eq:tau_broyden1}.

\begin{remark}
SSBroyden2 is the \emph{canonical} variant used in
Kiyani et al.\ (2022).  It was designed to avoid the numerical
instabilities that SSBroyden1 encounters when $b_k \approx 1$
(near-singular denominators in $\thk^{\text{SR1}}$).
\end{remark}

% ------------------------------------------------------------------
\subsection{Nocedal--Wright Initial Hessian Scaling}
\label{sec:nw_scaling}

When $\Hk = I$ (identity initialization), the first BFGS step is
essentially a steepest-descent step.  Nocedal and Wright
(\textit{Numerical Optimization}, Algorithm~6.1) suggest rescaling
$H_0$ after the first curvature pair:
\begin{equation}
\boxed{
  H_0 \gets \frac{\yk^\top \sk}{\yk^\top \yk}\, I
  = \frac{1}{\rho_k \,\|\yk\|^2}\, I
}
\label{eq:nw_scale}
\end{equation}
This sets the initial step length to match the observed curvature
magnitude, reducing the condition number of the first Hessian estimate.

% ------------------------------------------------------------------
\subsection{Power Transform for Hessian Conditioning}
\label{sec:power_transform}

Instead of minimizing $\mathcal{L}$ directly, we minimize
$\mathcal{L}_{\text{eff}} = \mathcal{L}^{1/p}$ for some $p \geq 1$.

\begin{theorem}[Condition Number Reduction]
If $\nabla^2 \mathcal{L}$ has condition number $\kappa$, then
$\nabla^2 \mathcal{L}^{1/p}$ has condition number approximately
$\kappa^{1/p}$ (in the regime where $\mathcal{L}$ is smooth and bounded
away from zero).
\end{theorem}

\textbf{Gradient of the transformed loss:}
\begin{equation}
  \nabla \mathcal{L}^{1/p} = \frac{1}{p}\,\mathcal{L}^{(1/p) - 1}\,\nabla \mathcal{L}.
  \label{eq:power_grad}
\end{equation}
When $\mathcal{L} \to 0$, the prefactor $\mathcal{L}^{(1/p)-1} \to \infty$
(for $p > 1$), which amplifies the gradient signal near convergence.

\begin{remark}
For $p = 2$: $\nabla \sqrt{\mathcal{L}} = \frac{\nabla \mathcal{L}}{2\sqrt{\mathcal{L}}}$.
The gradient is amplified by $1/(2\sqrt{\mathcal{L}})$, which diverges as
$\mathcal{L} \to 0$.  This can cause numerical instability in the final
stages of optimization.
\end{remark}

% ------------------------------------------------------------------
\subsection{Warm-Starting Across Resampling}
\label{sec:warmstart}

A distinctive feature of this PINN training setup is that the collocation
points are \textbf{resampled} every $N_{\text{change}}$ iterations.  When
the points change, the loss function $\mathcal{L}$ itself changes, so the
Hessian approximation $\Hk$ built on the old points is no longer exact.

Our implementation \textbf{warm-starts} $\Hk$: the inverse Hessian from
the end of one batch carries over as the initial estimate for the next.
A Cholesky factorization check is performed; if $\Hk$ has lost positive
definiteness (e.g., due to numerical drift), it is reset to $I$.

This warm-starting is only possible with dense BFGS methods.  L-BFGS-B
implicitly stores only recent curvature pairs, and SciPy's
implementation does not support warm-starting.


% ======================================================================
%  CHAPTER 5 — EXPERIMENT DESIGN
% ======================================================================
\section{Experiment Design}
\label{sec:design}

\subsection{Ablation Axes}

\begin{table}[H]
\centering
\caption{Summary of the 17 experiments.  All share the same network,
sampling, loss weights, and seed.  ``Canonical'' values are shown in bold.}
\label{tab:experiments}
\small
\begin{tabularx}{\textwidth}{llXllll}
\toprule
\textbf{Axis} & \textbf{Experiment} & \textbf{Variable} & \textbf{Adam} & \textbf{BFGS} & \textbf{Variant} & \textbf{Nchange} \\
\midrule
\multirow{2}{*}{Baseline}
  & canonical    & ---                    & \textbf{2000} & \textbf{10000} & \textbf{SSBroyden2} & \textbf{500} \\
  & adam\_only   & No BFGS                & 12000         & 0              & none                & ---          \\
\midrule
\multirow{5}{*}{A: Variant}
  & bfgs\_vanilla  & Standard BFGS        & 2000 & 10000 & BFGS\_scipy  & 500 \\
  & ssbfgs\_ol     & OL scaling           & 2000 & 10000 & SSBFGS\_OL   & 500 \\
  & ssbfgs\_ab     & AB scaling           & 2000 & 10000 & SSBFGS\_AB   & 500 \\
  & ssbroyden1     & Broyden $\phi=1$     & 2000 & 10000 & SSBroyden1   & 500 \\
  & lbfgs          & Limited memory       & 2000 & 10000 & L-BFGS-B     & 500 \\
\midrule
\multirow{3}{*}{B: Warmup}
  & warmup\_0      & No Adam              & 0    & 12000 & SSBroyden2 & 500 \\
  & warmup\_500    & Short warmup         & 500  & 11500 & SSBroyden2 & 500 \\
  & warmup\_5000   & Long warmup          & 5000 & 7000  & SSBroyden2 & 500 \\
\midrule
\multirow{2}{*}{C: Adam LR}
  & adam\_lr\_high  & $\text{lr} = 0.01$  & 2000 & 10000 & SSBroyden2 & 500 \\
  & adam\_lr\_low   & $\text{lr} = 10^{-4}$  & 2000 & 10000 & SSBroyden2 & 500 \\
\midrule
\multirow{2}{*}{D: Batch}
  & bfgs\_batch\_200  & Resample/200     & 2000 & 10000 & SSBroyden2 & 200 \\
  & bfgs\_batch\_1000 & Resample/1000    & 2000 & 10000 & SSBroyden2 & 1000 \\
\midrule
\multirow{2}{*}{E: Power}
  & power\_2       & $\sqrt{\mathcal{L}}$ & 2000 & 10000 & SSBroyden2 & 500 \\
  & power\_4       & $\mathcal{L}^{1/4}$  & 2000 & 10000 & SSBroyden2 & 500 \\
\midrule
F: Scale
  & hessian\_scaled & $H_0$ scaling       & 2000 & 10000 & SSBroyden2 & 500 \\
\bottomrule
\end{tabularx}
\end{table}

\subsection{Controlled Variables}

\begin{itemize}[nosep]
  \item \textbf{Network:} $[3, 20, 20, 20, 20, 1]$ with $\tanh$, variance scaling init, 1,361 params
  \item \textbf{Sampling:} $N_{\text{int}} = 10{,}000$, $N_0 = 1{,}000$, $N_b = 500$, RAD enabled
  \item \textbf{Loss weights:} $\lambda_{\text{pde}} = 1$, $\lambda_{\text{ic}} = 5$, $\lambda_{\text{bc}} = 5$
  \item \textbf{Seed:} 2 (PyTorch and NumPy)
  \item \textbf{Precision:} float64
  \item \textbf{Hardware:} OSCAR cluster (Brown CCV), 16 GB RAM, 4 CPU cores per job
\end{itemize}


% ======================================================================
%  CHAPTER 6 — RESULTS
% ======================================================================
\section{Results}
\label{sec:results}

\subsection{Overview}

\begin{table}[H]
\centering
\caption{Results for all 17 experiments, sorted by final loss (ascending).}
\label{tab:results}
\small
\begin{tabular}{llrrrr}
\toprule
\textbf{Experiment} & \textbf{Optimizer} & \textbf{Final Loss} & \textbf{Time (h)} & \textbf{Mass Drift} \\
\midrule
adam\_lr\_high      & SSBroyden2   & $2.36 \times 10^{-4}$ & 2.25 & 0.0112 \\
warmup\_5000        & SSBroyden2   & $2.86 \times 10^{-4}$ & 1.72 & 0.0027 \\
bfgs\_batch\_200    & SSBroyden2   & $3.30 \times 10^{-4}$ & 1.70 & 0.0025 \\
ssbfgs\_ab          & SSBFGS\_AB   & $3.40 \times 10^{-4}$ & 2.67 & 0.0094 \\
warmup\_0           & SSBroyden2   & $3.89 \times 10^{-4}$ & 2.13 & \textbf{0.0001} \\
adam\_lr\_low        & SSBroyden2   & $4.75 \times 10^{-4}$ & 2.24 & 0.0009 \\
warmup\_500         & SSBroyden2   & $6.99 \times 10^{-4}$ & 2.48 & 0.0134 \\
bfgs\_vanilla       & BFGS\_scipy  & $1.45 \times 10^{-3}$ & 2.14 & 0.0039 \\
ssbfgs\_ol          & SSBFGS\_OL   & $1.44 \times 10^{-3}$ & 2.42 & 0.0110 \\
power\_4            & SSBroyden2   & $3.82 \times 10^{-3}$ & 2.76 & 0.0044 \\
canonical           & SSBroyden2   & $5.10 \times 10^{-3}$ & 2.09 & 0.0172 \\
hessian\_scaled     & SSBroyden2   & $9.83 \times 10^{-3}$ & 3.51 & 0.0316 \\
lbfgs               & L-BFGS-B     & $1.05 \times 10^{-2}$ & 2.51 & 0.0099 \\
power\_2            & SSBroyden2   & $1.26 \times 10^{-2}$ & 2.77 & 0.0167 \\
adam\_only           & none (Adam)  & $2.72 \times 10^{-2}$ & 1.46 & 0.0090 \\
bfgs\_batch\_1000   & SSBroyden2   & $2.65 \times 10^{-2}$ & 1.70 & 0.0250 \\
ssbroyden1          & SSBroyden1   & $1.63$                & 2.94 & 0.0010 \\
\bottomrule
\end{tabular}
\end{table}

\begin{figure}[H]
  \centering
  \includegraphics[width=0.95\textwidth]{figures/final_loss_bar.pdf}
  \caption{Final loss for all 17 experiments (horizontal bar chart, log scale).
  The dashed blue line marks the canonical SSBroyden2 baseline.}
  \label{fig:final_loss_bar}
\end{figure}

\begin{figure}[H]
  \centering
  \includegraphics[width=0.95\textwidth]{figures/summary_heatmap.pdf}
  \caption{Normalized heatmap of final loss, total time, and mass drift
  for all experiments.  Green = better, red = worse.  Raw values shown
  in each cell.}
  \label{fig:heatmap}
\end{figure}

% ------------------------------------------------------------------
\subsection{Axis A: Optimizer Variant}

\begin{figure}[H]
  \centering
  \includegraphics[width=0.95\textwidth]{figures/optimizer_variant_comparison.pdf}
  \caption{Comparison of six optimizer variants (all with 2,000 Adam + 10,000 BFGS).
  Left: final loss (log scale).  Right: wall-clock time.}
  \label{fig:variant_comparison}
\end{figure}

\begin{figure}[H]
  \centering
  \includegraphics[width=0.95\textwidth]{figures/loss_curves_variants.pdf}
  \caption{Training loss curves for the six optimizer variants.
  Light traces show the Adam phase; solid traces show the BFGS phase.}
  \label{fig:variant_curves}
\end{figure}

\textbf{Key findings:}
\begin{itemize}[nosep]
  \item \textbf{SSBFGS-AB} achieves the best loss among pure variant comparisons
        ($3.40 \times 10^{-4}$), \textbf{15$\times$ better} than canonical SSBroyden2.
  \item \textbf{SSBroyden1 diverges} to loss $= 1.63$.  The DFP end of the Broyden
        family ($\thk = 1$) is catastrophically unsuitable for this problem.  DFP is known
        to be less robust than BFGS for general nonlinear optimization, and this confirms it.
  \item \textbf{Vanilla BFGS and SSBFGS-OL} perform comparably ($\sim 1.45 \times 10^{-3}$),
        both $\sim 3.5\times$ better than canonical.
  \item \textbf{L-BFGS-B underperforms} all dense methods ($1.05 \times 10^{-2}$).
        With only 1,361 parameters, the full Hessian is affordable (15 MB),
        and the limited-memory approximation loses essential curvature information.
        Additionally, L-BFGS-B cannot warm-start $\Hk$ across resampling events.
\end{itemize}

% ------------------------------------------------------------------
\subsection{Axis B: Adam Warmup Duration}

\begin{figure}[H]
  \centering
  \includegraphics[width=0.95\textwidth]{figures/warmup_sweep.pdf}
  \caption{Effect of Adam warmup duration on final loss (left) and mass
  conservation (right).  The relationship is U-shaped for loss.}
  \label{fig:warmup_sweep}
\end{figure}

\begin{figure}[H]
  \centering
  \includegraphics[width=0.95\textwidth]{figures/loss_curves_warmup.pdf}
  \caption{Training curves for the warmup sweep.  Dotted vertical lines
  mark the Adam$\to$BFGS transition for each configuration.}
  \label{fig:warmup_curves}
\end{figure}

\textbf{Key findings:}
\begin{itemize}[nosep]
  \item The \textbf{loss vs.\ warmup relationship is U-shaped}: the canonical
        2,000-epoch warmup is the \emph{worst} split in this sweep.
  \item \textbf{warmup\_0 (no Adam)} achieves $3.89 \times 10^{-4}$ and the
        \textbf{best mass conservation} of any experiment ($\Delta m = 0.0001$).
        This is a surprising result --- SSBroyden2 can train from random
        initialization without any first-order warmup.
  \item \textbf{warmup\_5000} achieves the lowest loss ($2.86 \times 10^{-4}$),
        suggesting that a well-trained starting point helps BFGS even with
        fewer remaining iterations.
  \item \textbf{Why is 2,000 worst?}  Hypothesis: 2,000 Adam epochs are enough to
        move the parameters into a region with complex Hessian structure, but not
        enough for Adam to find a clean basin.  BFGS then inherits a poor starting
        point.  Either very little warmup (BFGS handles everything) or extensive
        warmup (clean basin) works better.
\end{itemize}

% ------------------------------------------------------------------
\subsection{Axis C: Adam Learning Rate}

\begin{figure}[H]
  \centering
  \includegraphics[width=0.55\textwidth]{figures/adam_lr_sweep.pdf}
  \caption{Effect of Adam learning rate on final BFGS loss.
  Higher Adam LR produces a dramatically better handoff to BFGS.}
  \label{fig:lr_sweep}
\end{figure}

\textbf{Key findings:}
\begin{itemize}[nosep]
  \item \textbf{lr $= 0.01$ wins decisively} ($2.36 \times 10^{-4}$),
        achieving the best overall loss across all 17 experiments.
  \item Even $\text{lr} = 10^{-4}$ outperforms canonical $\text{lr} = 10^{-3}$,
        suggesting the default learning rate with the decay schedule
        overshoots the ideal Adam trajectory.
  \item The aggressive $\text{lr} = 0.01$ likely explores more of the landscape
        during the warmup phase, finding a basin that BFGS can efficiently
        exploit.
\end{itemize}

% ------------------------------------------------------------------
\subsection{Axis D: Collocation Resampling Frequency}

\begin{figure}[H]
  \centering
  \includegraphics[width=0.55\textwidth]{figures/batch_size_sweep.pdf}
  \caption{Effect of BFGS batch size ($N_{\text{change}}$) on final loss.
  More frequent resampling is dramatically better.}
  \label{fig:batch_sweep}
\end{figure}

\textbf{Key findings:}
\begin{itemize}[nosep]
  \item \textbf{batch\_size $= 200$} is \textbf{15$\times$ better} than canonical
        ($3.30 \times 10^{-4}$ vs.\ $5.10 \times 10^{-3}$).  This is one of
        the single most impactful hyperparameters.
  \item \textbf{batch\_size $= 1000$ essentially stalls} ($2.65 \times 10^{-2}$),
        barely improving on Adam.  The Hessian approximation built on stale
        collocation points becomes misleading after too many iterations.
  \item \textbf{Interpretation:}  When collocation points change, the loss surface
        changes.  Frequent resampling forces the optimizer to constantly re-explore,
        preventing overfitting to a fixed set of points.  The warm-started $\Hk$
        retains useful curvature information even after resampling, making the
        transition smooth.
\end{itemize}

% ------------------------------------------------------------------
\subsection{Axis E: Power Transform}

\begin{figure}[H]
  \centering
  \includegraphics[width=0.55\textwidth]{figures/power_transform.pdf}
  \caption{Power transform comparison.  The theoretical conditioning
  benefit does not materialize in practice.}
  \label{fig:power}
\end{figure}

\textbf{Key findings:}
\begin{itemize}[nosep]
  \item $p = 2$ (minimizing $\sqrt{\mathcal{L}}$) \textbf{hurts} performance
        ($1.26 \times 10^{-2}$, worse than canonical).  The gradient amplification
        near convergence ($\propto 1/\sqrt{\mathcal{L}}$) causes instability.
  \item $p = 4$ gives modest improvement ($3.82 \times 10^{-3}$) but is still
        dominated by other interventions (variant choice, batch size, LR).
  \item The theoretical $\kappa^{1/p}$ conditioning benefit does not outweigh
        the numerical difficulties introduced by the gradient rescaling.
\end{itemize}

% ------------------------------------------------------------------
\subsection{Axis F: Initial Hessian Scaling}

\begin{figure}[H]
  \centering
  \includegraphics[width=0.7\textwidth]{figures/hessian_scaling.pdf}
  \caption{Nocedal--Wright $H_0$ scaling: increases both loss and wall time.}
  \label{fig:hessian}
\end{figure}

\textbf{Key findings:}
\begin{itemize}[nosep]
  \item Nocedal--Wright scaling \textbf{hurts}: $9.83 \times 10^{-3}$ vs.\
        $5.10 \times 10^{-3}$ (canonical), and $3.51$~h vs.\ $2.09$~h.
  \item The one-time $H_0 = (\yk^\top \sk / \yk^\top \yk) I$ rescaling
        likely sets a bad initial scale that takes many iterations to correct,
        especially with the warm-start/resample scheme where the first
        batch's curvature may not be representative.
\end{itemize}

% ------------------------------------------------------------------
\subsection{Loss vs.\ Compute (Pareto Front)}

\begin{figure}[H]
  \centering
  \includegraphics[width=0.75\textwidth]{figures/pareto_loss_time.pdf}
  \caption{Scatter plot of final loss vs.\ total wall time.  Points
  in the lower-left corner are Pareto-optimal (low loss, fast).}
  \label{fig:pareto}
\end{figure}

The Pareto-optimal configurations are:
\begin{enumerate}[nosep]
  \item \textbf{warmup\_5000} ($2.86 \times 10^{-4}$, 1.72 h) --- best loss/time ratio
  \item \textbf{bfgs\_batch\_200} ($3.30 \times 10^{-4}$, 1.70 h) --- equally fast, slightly higher loss
  \item \textbf{adam\_lr\_high} ($2.36 \times 10^{-4}$, 2.25 h) --- best absolute loss
\end{enumerate}

% ------------------------------------------------------------------
\subsection{Mass Conservation}

\begin{figure}[H]
  \centering
  \includegraphics[width=0.95\textwidth]{figures/mass_drift_bar.pdf}
  \caption{Mass drift for all experiments.  warmup\_0 achieves near-zero
  drift, suggesting that skipping Adam preserves physical conservation
  laws better.}
  \label{fig:mass}
\end{figure}

\textbf{Key findings:}
\begin{itemize}[nosep]
  \item \textbf{warmup\_0} has the best mass conservation ($\Delta m = 0.0001$),
        two orders of magnitude better than most experiments.
  \item Adam-only shows moderate drift ($0.009$) despite no BFGS phase.
  \item \textbf{hessian\_scaled} has the worst mass conservation ($0.032$),
        consistent with its poor loss.
  \item There is a weak negative correlation between final loss and mass
        drift (better loss $\not\Rightarrow$ better conservation), suggesting
        these are partially independent objectives.
\end{itemize}


% ======================================================================
%  CHAPTER 7 — DISCUSSION
% ======================================================================
\section{Discussion}
\label{sec:discussion}

\subsection{Why Does SSBFGS-AB Win Among Variants?}

The Al-Baali scaling $\tp = \min(1, 1/b_k)$ is a \emph{conservative
one-sided clamp}: it never shrinks the inverse Hessian (which would make
steps more conservative), but freely expands it when the curvature
estimate is too large.  In the PINN context, where the loss landscape
changes due to collocation resampling, this asymmetry is beneficial:
\begin{itemize}[nosep]
  \item After resampling, the true curvature may be smaller than what
        $\Hk$ encodes.  AB expands $\Hk$ to compensate.
  \item If the curvature grows, AB leaves $\Hk$ unchanged (safe default)
        while OL would aggressively shrink it.
\end{itemize}

\subsection{Why Does SSBroyden1 Diverge?}

SSBroyden1 uses $\thk$ choices that favor the DFP end of the Broyden
family.  DFP is known to be sensitive to inaccurate line searches
(it requires \emph{exact} line searches for theoretical guarantees).
With approximate Wolfe line searches and the added complication of
collocation resampling, the DFP-biased updates accumulate error in
$\Hk$ until it becomes a poor approximation, leading to divergence.

\subsection{The Resampling Frequency Insight}

The $80\times$ gap between batch\_200 and batch\_1000 is one of the
most striking results.  This suggests a fundamental tension in PINN
training:
\begin{itemize}[nosep]
  \item \textbf{Too rare resampling} ($N_{\text{change}} = 1000$):
        BFGS overfits to the current collocation points.  The loss
        decreases on those points but does not generalize.
  \item \textbf{Too frequent resampling} would waste curvature information
        built up in $\Hk$, though we did not reach this regime.
  \item The \textbf{sweet spot} is frequent enough to prevent overfitting
        but rare enough for BFGS to make meaningful progress.
\end{itemize}

\subsection{Why Can SSBroyden2 Train from Scratch?}

The success of warmup\_0 challenges the conventional wisdom that
PINNs need first-order warmup.  SSBroyden2's self-scaling mechanism
automatically adjusts the step size based on observed curvature,
effectively performing its own ``warmup''.  The Wolfe line search
provides safeguards against too-large steps.  The result: quasi-Newton
methods are more self-sufficient than previously assumed, at least
for small networks.


% ======================================================================
%  CHAPTER 8 — NEXT STEPS
% ======================================================================
\section{Next Steps}
\label{sec:next}

Based on the experimental results and analysis, we propose the following
directions for future work:

\subsection{Immediate Follow-ups}

\begin{enumerate}
  \item \textbf{Combine the winners.}  Run a ``best-of'' configuration
        combining SSBFGS-AB (best variant) + batch\_size $= 200$ (best
        resampling) + Adam lr $= 0.01$ (best LR) + 5,000 warmup (best
        warmup for loss).  These are independent axes that should compound.

  \item \textbf{Fine-grained resampling sweep.}  Test $N_{\text{change}} \in
        \{50, 100, 150, 200, 300\}$ to find the optimal frequency and
        determine whether the improvement continues below 200.

  \item \textbf{Multi-seed robustness.}  All 17 experiments used seed $= 2$.
        Re-run the top 5 configurations with seeds $\{1, 2, 3, 4, 5\}$ to
        obtain confidence intervals and verify that the rankings are stable
        under different random initializations.

  \item \textbf{Investigate warmup\_0's mass conservation.}  The near-zero
        mass drift of the no-warmup configuration is remarkable.
        Investigate whether Adam introduces a systematic bias that violates
        the conservation law, and whether this finding generalizes to other
        conserved-quantity PDEs.
\end{enumerate}

\subsection{Scaling and Generalization}

\begin{enumerate}[resume]
  \item \textbf{Larger networks.}  Test with $[3, 40, 40, 40, 40, 1]$
        ($\sim 5{,}000$ params) and $[3, 64, 64, 64, 64, 1]$ ($\sim 12{,}000$
        params).  At some point the $O(n^2)$ storage of dense BFGS becomes
        prohibitive and L-BFGS-B may become competitive.

  \item \textbf{GPU acceleration.}  Port the Hessian update to PyTorch
        (GPU) to eliminate the CPU$\leftrightarrow$GPU transfer bottleneck.
        This would enable testing on larger networks where the current
        NumPy-based Hessian computation is the time bottleneck.

  \item \textbf{Other PDEs.}  Apply the winning configuration to:
    \begin{itemize}[nosep]
      \item Burgers' equation (1D, simpler baseline)
      \item Allen--Cahn equation (related phase-field model, no biharmonic)
      \item Navier--Stokes equations (higher-dimensional, more challenging)
    \end{itemize}

  \item \textbf{Adaptive resampling frequency.}  Instead of fixed
        $N_{\text{change}}$, adapt the frequency based on a staleness
        metric (e.g., the relative change in $\Hk$ eigenvalues or the
        loss on a held-out validation set of collocation points).
\end{enumerate}

\subsection{Theoretical Directions}

\begin{enumerate}[resume]
  \item \textbf{Convergence analysis.}  Develop convergence rate bounds
        for SSBFGS-AB in the non-stationary setting (changing loss function
        due to resampling).  Existing quasi-Newton convergence theory assumes
        a fixed objective.

  \item \textbf{Condition number tracking.}  Log $\kappa(\Hk)$ over training
        to understand how different variants control the Hessian condition
        number, and correlate with convergence speed.

  \item \textbf{Hybrid strategies.}  Explore switching between SSBFGS-AB
        and SSBroyden2 during training based on the value of $b_k$
        (using AB when $b_k > 1$ and SSBroyden2 otherwise).
\end{enumerate}


% ======================================================================
%  REFERENCES
% ======================================================================
\newpage
\section*{References}

\begin{enumerate}[label={[\arabic*]},leftmargin=*,nosep]
  \item Kiyani, E., Silber, S., Shapira, Y., and Karniadakis, G.E.
        ``Optimizing the Optimizer for Physics-Informed Neural Networks.''
        \textit{Physical Review E}, 106(6):065303, 2022.

  \item Nocedal, J.\ and Wright, S.J.
        \textit{Numerical Optimization}, 2nd ed.
        Springer, 2006.

  \item Oren, S.S.\ and Luenberger, D.G.
        ``Self-Scaling Variable Metric (SSVM) Algorithms. Part I: Criteria
        and Sufficient Conditions for Scaling a Class of Algorithms.''
        \textit{Management Science}, 20(5):845--862, 1974.

  \item Al-Baali, M.
        ``Numerical Experience with a Class of Self-Scaling Quasi-Newton
        Algorithms.''
        \textit{Journal of Optimization Theory and Applications},
        96(3):533--553, 1998.

  \item Raissi, M., Perdikaris, P., and Karniadakis, G.E.
        ``Physics-Informed Neural Networks: A Deep Learning Framework for
        Solving Forward and Inverse Problems Involving Nonlinear Partial
        Differential Equations.''
        \textit{Journal of Computational Physics}, 378:686--707, 2019.

  \item Cahn, J.W.\ and Hilliard, J.E.
        ``Free Energy of a Nonuniform System. I.\ Interfacial Free Energy.''
        \textit{The Journal of Chemical Physics}, 28(2):258--267, 1958.

  \item Liu, D.C.\ and Nocedal, J.
        ``On the Limited Memory BFGS Method for Large Scale Optimization.''
        \textit{Mathematical Programming}, 45(1):503--528, 1989.
\end{enumerate}


\end{document}
